% =========================================================
% GUIDE : CAPTURES D'ÉCRAN À INSÉRER — PaieZone RH
% Rapport PFE — Méthodologie Scrum — 4 sprints
% =========================================================
% Pour chaque interface :
%   - Sous-section exacte (label LaTeX)
%   - Nom de fichier image suggéré
%   - Référence figure déjà déclarée dans le rapport
%   - Ce qui doit apparaître à l'écran
% =========================================================
% IMPORTANT : Remplace \includegraphics[...]{scrum.png}
% par \includegraphics[...]{NOM_FICHIER_IMAGE} dans ton
% fichier .tex correspondant.
% =========================================================

% =========================================================
% CHAPITRE 3 — Sprint 1 : Infrastructure & Admin
% Objectif : Socle technique, sécurité, multi-tenant
% =========================================================

% ---------------------------------------------------------
% Sous-section : \subsection{Inscription et initialisation
%                de l'espace société}
% Position : Après le texte introductif, avant la légende
% Label figure : fig:ui-inscription
% Fichier image : sp1_inscription.png
% ---------------------------------------------------------
% Ce que doit montrer la capture :
%   Page d'inscription de PaieZone.
%   Deux blocs de formulaire visibles simultanément :
%   (1) Informations personnelles de l'admin
%       (prénom, nom, e-mail, mot de passe)
%   (2) Informations légales de l'entreprise
%       (raison sociale, matricule fiscal, adresse)
%   Bouton "S'inscrire" en bas.
%   URL : /register ou /signup
%
% Code LaTeX (remplace scrum.png) :
% \includegraphics[width=0.7\textwidth]{sp1_inscription}

% ---------------------------------------------------------
% Sous-section : \subsection{Authentification et accès}
% Capture 1/5 — Page de connexion
% Label figure : fig:ui-login
% Fichier image : sp1_login.png
% ---------------------------------------------------------
% Ce que doit montrer la capture :
%   Page de connexion.
%   Champ e-mail + champ mot de passe + bouton "Se connecter".
%   Logo PaieZone visible en haut.
%   URL : /login
%
% Code LaTeX :
% \includegraphics[width=0.75\textwidth]{sp1_login}

% ---------------------------------------------------------
% Sous-section : \subsection{Authentification et accès}
% Capture 2/5 — Dashboard Super Admin
% Label figure : fig:ui-dashboard-Superadmin
% Fichier image : sp1_dashboard_superadmin.png
% ---------------------------------------------------------
% Ce que doit montrer la capture :
%   Tableau de bord Super Admin après connexion.
%   Widgets statistiques visibles :
%   - Nombre total d'entreprises clientes
%   - Abonnements actifs / expirés
%   - Graphique ou liste des sociétés
%   Barre de navigation latérale avec les modules SA
%   URL : /saas-dash ou /super-admin/dashboard
%
% Code LaTeX :
% \includegraphics[width=0.75\textwidth]{sp1_dashboard_superadmin}

% ---------------------------------------------------------
% Sous-section : \subsection{Authentification et accès}
% Capture 3/5 — Dashboard Administrateur
% Label figure : fig:ui-dashboard-admin
% Fichier image : sp1_dashboard_admin.png
% ---------------------------------------------------------
% Ce que doit montrer la capture :
%   Tableau de bord Admin Entreprise.
%   Widgets : effectif total, utilisateurs actifs,
%   abonnement en cours (plan + expiration),
%   dernières actions du journal d'audit.
%   Barre de navigation avec modules Admin.
%   URL : /admin/dashboard
%
% Code LaTeX :
% \includegraphics[width=0.75\textwidth]{sp1_dashboard_admin}

% ---------------------------------------------------------
% Sous-section : \subsection{Authentification et accès}
% Capture 4/5 — Dashboard Employé
% Label figure : fig:ui-dashboard-employe
% Fichier image : sp1_dashboard_employe.png
% ---------------------------------------------------------
% Ce que doit montrer la capture :
%   Tableau de bord Employé.
%   Widgets : solde de congés restants, prochain bulletin,
%   dernières demandes en cours.
%   Navigation limitée aux modules Employé.
%   URL : /employee/dashboard
%
% Code LaTeX :
% \includegraphics[width=0.75\textwidth]{sp1_dashboard_employe}

% ---------------------------------------------------------
% Sous-section : \subsection{Authentification et accès}
% Capture 5/5 — Dashboard RH Comptable
% Label figure : fig:ui-dashboard-rh
% Fichier image : sp1_dashboard_rh.png
% ---------------------------------------------------------
% Ce que doit montrer la capture :
%   Tableau de bord RH Comptable.
%   Widgets : effectif actif, congés en attente de validation,
%   bulletins à calculer pour la période courante.
%   Navigation RH (Employés, Paie, Congés, Chatbot).
%   URL : /rh/dashboard
%
% Code LaTeX :
% \includegraphics[width=0.75\textwidth]{sp1_dashboard_rh}

% ---------------------------------------------------------
% Sous-section : \subsection{Gestion des utilisateurs}
% Capture 1/2 — Liste des utilisateurs
% Label figure : fig:ui-users-list
% Fichier image : sp1_users_list.png
% ---------------------------------------------------------
% Ce que doit montrer la capture :
%   Tableau paginé des utilisateurs de la société.
%   Colonnes : Nom, Prénom, E-mail, Rôle, Statut (Actif/Inactif).
%   Boutons : Ajouter, Inviter, Modifier, Désactiver.
%   Barre de recherche en haut.
%   URL : /admin/users
%
% Code LaTeX :
% \includegraphics[width=0.70\textwidth]{sp1_users_list}

% ---------------------------------------------------------
% Sous-section : \subsection{Gestion des utilisateurs}
% Capture 2/2 — Formulaire création/modification
% Label figure : fig:ui-users-form
% Fichier image : sp1_users_form.png
% ---------------------------------------------------------
% Ce que doit montrer la capture :
%   Modal ou page de formulaire.
%   Champs : Prénom, Nom, E-mail, Rôle (liste déroulante).
%   Deux boutons en bas : "Annuler" et "Enregistrer".
%   Si formulaire d'invitation : champ e-mail uniquement
%   + message "Un lien d'activation sera envoyé".
%
% Code LaTeX :
% \includegraphics[width=0.70\textwidth]{sp1_users_form}

% ---------------------------------------------------------
% Sous-section : \subsection{Journal d'audit}
% Label figure : fig:ui-audit-list
% Fichier image : sp1_audit_list.png
% ---------------------------------------------------------
% Ce que doit montrer la capture :
%   Tableau du journal d'audit.
%   Colonnes : Date/Heure, Utilisateur, Action (CREATE/UPDATE/DELETE),
%   Entité cible, Valeur avant, Valeur après.
%   Filtres visibles : par date, par utilisateur, par type d'action.
%   URL : /admin/audit-logs
%
% Code LaTeX :
% \includegraphics[width=0.70\textwidth]{sp1_audit_list}

% ---------------------------------------------------------
% Sous-section : \subsection{Gestion Organisationnelle}
% Label figure : fig:ui-org
% Fichier image : sp1_org_structure.png
% ---------------------------------------------------------
% Ce que doit montrer la capture :
%   Interface de gestion des départements et postes.
%   Liste des départements avec bouton "Ajouter département".
%   Pour chaque département : liste des postes associés
%   + bouton "Ajouter poste".
%   Statut Actif/Inactif visible pour chaque ligne.
%   URL : /admin/organization
%
% Code LaTeX :
% \includegraphics[width=0.70\textwidth]{sp1_org_structure}


% =========================================================
% CHAPITRE 4 — Sprint 2 : Gestion du Capital Humain
% Objectif : Employés, contrats, documents administratifs
% =========================================================

% ---------------------------------------------------------
% Sous-section : \subsection{Gestion des employés}
% Capture 1/3 — Liste des employés
% Label figure : fig:ui-employees-list
% Fichier image : sp2_employees_list.png
% ---------------------------------------------------------
% Ce que doit montrer la capture :
%   Tableau paginé du personnel.
%   Colonnes : Matricule, Nom complet, Département, Poste, Statut.
%   Barre de recherche multicritère + filtres Département/Statut.
%   Boutons : Ajouter, Voir profil, Modifier, Désactiver.
%   URL : /rh/employees
%
% Code LaTeX :
% \includegraphics[width=0.70\textwidth]{sp2_employees_list}

% ---------------------------------------------------------
% Sous-section : \subsection{Gestion des employés}
% Capture 2/3 — Formulaire de création/modification
% Label figure : fig:ui-employee-form
% Fichier image : sp2_employee_form.png
% ---------------------------------------------------------
% Ce que doit montrer la capture :
%   Formulaire multi-onglets ou multi-sections.
%   Visible : Identité (nom FR + arabe, CIN, CNSS, genre),
%   Fiscal (situation familiale, enfants, chef famille),
%   Contact (adresse, e-mail pro, téléphone, RIB),
%   Poste (département, poste, date embauche).
%   Champs obligatoires marqués (*).
%   URL : /rh/employees/new ou /rh/employees/{id}/edit
%
% Code LaTeX :
% \includegraphics[width=0.70\textwidth]{sp2_employee_form}

% ---------------------------------------------------------
% Sous-section : \subsection{Gestion des employés}
% Capture 3/3 — Vue détaillée du profil
% Label figure : fig:ui-employee-profile
% Fichier image : sp2_employee_profile.png
% ---------------------------------------------------------
% Ce que doit montrer la capture :
%   Page de profil complet d'un employé.
%   Onglets visibles : Informations, Contrats, Documents, Historique.
%   En-tête : photo (optionnelle), matricule, nom, poste, département.
%   Onglet "Informations" actif montrant les données RH.
%   URL : /rh/employees/{id}
%
% Code LaTeX :
% \includegraphics[width=0.70\textwidth]{sp2_employee_profile}

% ---------------------------------------------------------
% Sous-section : \subsection{Gestion des contrats de travail}
% Label figure : fig:ui-contracts-list
% Fichier image : sp2_contracts_section.png
% ---------------------------------------------------------
% Ce que doit montrer la capture :
%   Onglet "Contrats" du profil employé.
%   Tableau : Référence, Type (CDI/CDD/CIVP/KARAMA/Intérimaire/Stage),
%   Date début, Date fin, Salaire brut, Statut (ACTIF/RÉSILIÉ).
%   Bouton "Associer un contrat" visible.
%   Pour le contrat actif : boutons "Modifier" et "Résilier".
%   URL : /rh/employees/{id} (onglet Contrats)
%
% Code LaTeX :
% \includegraphics[width=0.70\textwidth]{sp2_contracts_section}

% ---------------------------------------------------------
% Sous-section : \subsection{Gestion des documents administratifs}
% Label figure : fig:ui-documents
% Fichier image : sp2_documents_section.png
% ---------------------------------------------------------
% Ce que doit montrer la capture :
%   Onglet "Documents" du profil employé.
%   Tableau : Nom fichier, Type (CONTRAT/ATTESTATION/DIPLOME/AUTRE),
%   Date de dépôt, Actions (Télécharger, Supprimer).
%   Bouton "Ajouter un document" avec modal de sélection de type.
%   URL : /rh/employees/{id} (onglet Documents)
%
% Code LaTeX :
% \includegraphics[width=0.70\textwidth]{sp2_documents_section}


% =========================================================
% CHAPITRE 5 — Sprint 3 : Paie, Congés, Avances
% Objectif : Moteur de calcul, bulletins PDF, LF 2026
% =========================================================

% ---------------------------------------------------------
% Sous-section : \subsection{Gestion des paramètres réglementaires}
% Label figure : fig:ui-params
% Fichier image : sp3_params_reglementaires.png
% ---------------------------------------------------------
% Ce que doit montrer la capture :
%   Interface Super Admin — Paramètres légaux.
%   Tableau des constantes : Taux CNSS salarié (9,18% ou 9,68%),
%   CNSS patronal, CAVIS (1%), CSS (0,5%), TFP (1%),
%   Barème IRPP (tranches), SMIG.
%   Boutons "Ajouter paramètre" et "Modifier".
%   Colonne "Année d'application" (ex : 2026).
%   URL : /super-admin/payroll-params
%
% Code LaTeX :
% \includegraphics[width=0.70\textwidth]{sp3_params_reglementaires}

% ---------------------------------------------------------
% Sous-section : \subsection{Gestion des périodes de paie}
% Label figure : fig:ui-periodes
% Fichier image : sp3_periodes_paie.png
% ---------------------------------------------------------
% Ce que doit montrer la capture :
%   Interface des périodes de paie.
%   Tableau : Mois/Année, Statut (OUVERTE / CALCULÉE / CLÔTURÉE),
%   Date de création, Nombre de bulletins.
%   Bouton "Créer nouvelle période" visible.
%   Actions par ligne : "Calculer" (si OUVERTE), "Clôturer" (si CALCULÉE).
%   URL : /rh/payroll-periods
%
% Code LaTeX :
% \includegraphics[width=0.70\textwidth]{sp3_periodes_paie}

% ---------------------------------------------------------
% Sous-section : \subsection{Calcul et visualisation des bulletins de paie}
% Label figure : fig:ui-bulletins
% Fichier image : sp3_bulletins_paie.png
% ---------------------------------------------------------
% Ce que doit montrer la capture :
%   Interface des bulletins d'une période.
%   Tableau : Matricule, Employé, Salaire brut, CNSS salarié,
%   CAVIS, CSS, IRPP, Net à payer.
%   Bouton "Calculer la période" en haut.
%   Pour chaque bulletin : bouton "Télécharger PDF" et "Envoyer par e-mail".
%   OU : aperçu d'un bulletin PDF généré (pages avec mentions légales).
%   URL : /rh/payroll-periods/{id}/payslips
%
% Code LaTeX :
% \includegraphics[width=0.70\textwidth]{sp3_bulletins_paie}

% ---------------------------------------------------------
% ASTUCE : Ajoute une 2e capture — aperçu du bulletin PDF généré
% Label figure : fig:ui-bulletin-pdf  (NOUVELLE figure à ajouter)
% Fichier image : sp3_bulletin_pdf.png
% ---------------------------------------------------------
% Ce que doit montrer la capture :
%   Aperçu d'un bulletin de paie généré en PDF (fenêtre navigateur).
%   En-tête : Logo société + Période (ex: Mai 2026).
%   Corps : Rubriques de paie (Salaire brut, Primes, CNSS, IRPP, Net).
%   Pied de page : Signature + Cachet.
%   (Ce PDF est produit par JasperReports / Flying Saucer)
%
% Code LaTeX à ajouter après fig:ui-bulletins :
% \begin{figure}[H]
%   \centering
%   \includegraphics[width=0.70\textwidth]{sp3_bulletin_pdf}
%   \caption{Aperçu du bulletin de paie généré au format PDF}
%   \label{fig:ui-bulletin-pdf}
% \end{figure}

% ---------------------------------------------------------
% Sous-section : \subsection{Gestion des primes et des avances}
% Label figure : fig:ui-primes
% Fichier image : sp3_primes_avances.png
% ---------------------------------------------------------
% Ce que doit montrer la capture :
%   Interface de gestion des éléments variables.
%   Deux onglets : "Primes" et "Avances".
%   Onglet Primes : tableau des primes ajoutées pour la période
%   (Employé, Type prime, Montant, Statut).
%   OU onglet Avances : tableau des demandes
%   (Employé, Montant demandé, Date, Statut REQUESTED/APPROVED/REJECTED),
%   + boutons "Approuver" et "Refuser".
%   URL : /rh/payroll-periods/{id}/variables
%
% Code LaTeX :
% \includegraphics[width=0.70\textwidth]{sp3_primes_avances}

% ---------------------------------------------------------
% Sous-section : \subsection{Gestion des congés et des absences}
% Label figure : fig:ui-conges
% Fichier image : sp3_conges.png
% ---------------------------------------------------------
% Ce que doit montrer la capture :
%   Interface de gestion des congés (vue RH ou Employé).
%
%   VUE RH (recommandée) :
%   Tableau des demandes en attente : Employé, Dates, Type,
%   Jours demandés, Statut (EN ATTENTE / APPROUVÉ / REFUSÉ).
%   Boutons "Valider" et "Refuser" par ligne.
%
%   VUE EMPLOYÉ (alternative) :
%   Calendrier interactif avec les congés posés.
%   Compteur de solde en temps réel (ex: "12 jours restants").
%   Bouton "Soumettre une demande".
%
%   URL : /rh/leave-requests ou /employee/leaves
%
% Code LaTeX :
% \includegraphics[width=0.70\textwidth]{sp3_conges}


% =========================================================
% CHAPITRE 6 — Sprint 4 : Comptabilité, Conformité, Chatbot
% Objectif : Export SAGE, déclarations légales, PaieBot IA
% =========================================================

% ---------------------------------------------------------
% Sous-section : \subsection{Gestion comptable et export SAGE}
% Label figure : fig:ui-comptabilite
% Fichier image : sp4_comptabilite_sage.png
% ---------------------------------------------------------
% Ce que doit montrer la capture :
%   Interface de l'onglet Comptabilité.
%   Tableau du plan comptable : Numéro de compte, Libellé,
%   Type (DEBIT/CREDIT), Actions (Modifier, Supprimer).
%   Bouton "Exporter SAGE (.csv)" visible en haut à droite.
%   URL : /rh/accounting ou /admin/accounting
%
% Code LaTeX :
% \includegraphics[width=0.70\textwidth]{sp4_comptabilite_sage}

% ---------------------------------------------------------
% Sous-section : \subsection{Conformité réglementaire}
% Label figure : fig:ui-conformite
% Fichier image : sp4_conformite_declarations.png
% ---------------------------------------------------------
% Ce que doit montrer la capture :
%   Interface de l'onglet Conformité.
%   Quatre boutons/cards de génération :
%   (1) Déclaration CNSS salarié (trimestrielle)
%   (2) Déclaration CNSS patronal (trimestrielle)
%   (3) Déclaration CAVIS (trimestrielle)
%   (4) Déclaration IRPP (annuelle)
%   Chaque card : période sélectionnable + bouton "Générer".
%   URL : /rh/compliance
%
% Code LaTeX :
% \includegraphics[width=0.70\textwidth]{sp4_conformite_declarations}

% ---------------------------------------------------------
% Sous-section : \subsection{Documents administratifs RH}
% Label figure : fig:ui-docs-rh
% Fichier image : sp4_docs_rh_generation.png
% ---------------------------------------------------------
% Ce que doit montrer la capture :
%   Interface de génération des documents officiels.
%   Accessible depuis la liste des employés ou le profil.
%   Deux boutons visibles :
%   - "Attestation de salaire (PDF)"
%   - "Certificat de travail (PDF)"
%   OU : aperçu d'une attestation PDF générée (template légal).
%   URL : /rh/employees (boutons d'action sur la ligne)
%
% Code LaTeX :
% \includegraphics[width=0.70\textwidth]{sp4_docs_rh_generation}

% ---------------------------------------------------------
% Sous-section : \subsection{Assistant conversationnel RH (PaieBot)}
% Label figure : fig:ui-chatbot
% Fichier image : sp4_paibot_chatbot.png
% ---------------------------------------------------------
% Ce que doit montrer la capture :
%   Composant chatbot flottant en bas à droite de l'écran.
%   Bulle de chat ouverte avec une conversation exemple.
%   Questions et réponses visibles (en français) :
%   Ex Employé : "Quel est mon salaire net ce mois ?"
%   Réponse PaieBot : "Votre salaire net pour Mai 2026 est de 2 450 TND."
%   Ex Admin : "Combien d'employés actifs dans la société ?"
%   Réponse : "Votre société compte 47 employés actifs."
%   Icône PaieBot visible + barre de saisie en bas.
%   Contexte de la page arrière-plan : Dashboard Admin ou RH.
%
% Code LaTeX :
% \includegraphics[width=0.70\textwidth]{sp4_paibot_chatbot}

% ---------------------------------------------------------
% ASTUCE : Ajoute une 2e capture — réponse RAG réglementaire
% Label figure : fig:ui-chatbot-rag  (NOUVELLE figure à ajouter)
% Fichier image : sp4_paibot_rag.png
% ---------------------------------------------------------
% Ce que doit montrer la capture :
%   Même composant chatbot.
%   Question réglementaire visible :
%   Ex : "Quel est le taux CNSS salarié en 2026 ?"
%   Réponse PaieBot issue du pipeline LEGAL_RAG :
%   "Selon la Loi de Finances 2026, le taux CNSS salarié est de 9,68%..."
%   Source documentaire citée en bas de la réponse (optionnel).
%   Badge "Réglementation" ou indication du type de pipeline visible.
%
% Code LaTeX à ajouter après fig:ui-chatbot :
% \begin{figure}[H]
%   \centering
%   \includegraphics[width=0.70\textwidth]{sp4_paibot_rag}
%   \caption{PaieBot --- Réponse réglementaire via pipeline RAG (LEGAL\_RAG)}
%   \label{fig:ui-chatbot-rag}
% \end{figure}


% =========================================================
% RÉCAPITULATIF GÉNÉRAL — Toutes les captures à prendre
% =========================================================
%
% SPRINT 1 (Chapitre 3) — 9 captures obligatoires + 0 optionnelles :
%   1. sp1_inscription.png         → fig:ui-inscription
%   2. sp1_login.png               → fig:ui-login
%   3. sp1_dashboard_superadmin.png→ fig:ui-dashboard-Superadmin
%   4. sp1_dashboard_admin.png     → fig:ui-dashboard-admin
%   5. sp1_dashboard_employe.png   → fig:ui-dashboard-employe
%   6. sp1_dashboard_rh.png        → fig:ui-dashboard-rh
%   7. sp1_users_list.png          → fig:ui-users-list
%   8. sp1_users_form.png          → fig:ui-users-form
%   9. sp1_audit_list.png          → fig:ui-audit-list
%  10. sp1_org_structure.png       → fig:ui-org
%
% SPRINT 2 (Chapitre 4) — 5 captures obligatoires :
%   1. sp2_employees_list.png      → fig:ui-employees-list
%   2. sp2_employee_form.png       → fig:ui-employee-form
%   3. sp2_employee_profile.png    → fig:ui-employee-profile
%   4. sp2_contracts_section.png   → fig:ui-contracts-list
%   5. sp2_documents_section.png   → fig:ui-documents
%
% SPRINT 3 (Chapitre 5) — 5 captures obligatoires + 1 recommandée :
%   1. sp3_params_reglementaires.png → fig:ui-params
%   2. sp3_periodes_paie.png         → fig:ui-periodes
%   3. sp3_bulletins_paie.png        → fig:ui-bulletins
%   4. sp3_primes_avances.png        → fig:ui-primes
%   5. sp3_conges.png                → fig:ui-conges
%   6. sp3_bulletin_pdf.png          → fig:ui-bulletin-pdf [RECOMMANDÉE]
%
% SPRINT 4 (Chapitre 6) — 4 captures obligatoires + 1 recommandée :
%   1. sp4_comptabilite_sage.png         → fig:ui-comptabilite
%   2. sp4_conformite_declarations.png   → fig:ui-conformite
%   3. sp4_docs_rh_generation.png        → fig:ui-docs-rh
%   4. sp4_paibot_chatbot.png            → fig:ui-chatbot
%   5. sp4_paibot_rag.png                → fig:ui-chatbot-rag [RECOMMANDÉE]
%
% TOTAL : 24 captures obligatoires + 2 recommandées = 26 images
%
% =========================================================
% COMMENT REMPLACER scrum.png DANS TES .tex :
% =========================================================
% Dans chaque chapitre (chapitre3.tex ... chapitre6.tex),
% remplace chaque occurrence de :
%   \includegraphics[...]{scrum.png}
% par le nom de fichier image correspondant indiqué ci-dessus.
%
% Place tes captures dans le même dossier que tes .tex
% (C:\Users\chayma\PaieZoneRH\pfe_rapport_extracted\pfe_rapport\)
%
% Formats acceptés par LaTeX : .png, .jpg, .pdf
% Résolution recommandée : 1280x720 minimum (format 16:9)
% =========================================================
